\chapter{Conclusioni Generali} \label{chap:conclusioni}

Il progetto ha affrontato e risolto il problema del controllo di un robot uniciclo in condizioni operative degradate, sviluppando un'architettura completa che integra modellistica, controllo, diagnosi e riconfigurazione.

\section{Sintesi dei Risultati}
Il lavoro svolto ha permesso di raggiungere i seguenti obiettivi specifici:

\begin{itemize}
    \item \textbf{Validazione del Modello (Cap. \ref{chap:modello}):} La simulazione in catena aperta ha confermato la correttezza del modello cinematico affine, ponendo basi solide per le fasi successive.
    
    \item \textbf{Precisione del Controllo (Cap. \ref{chap:tracking}):} L'adozione della tecnica di \textit{Feedback Linearization} ha garantito prestazioni di tracking eccellenti in condizioni nominali. L'errore di posizione sulla traiettoria circolare è risultato trascurabile (\( < 2 \) cm), dimostrando la superiorità di questo approccio rispetto a controllori lineari classici per sistemi non-olonomi.
    
    \item \textbf{Impatto dei Guasti (Cap. \ref{chap:guasti}):} Le simulazioni in presenza di guasti hanno evidenziato la vulnerabilità del controllo nominale. È stato dimostrato che un guasto asimmetrico (es. FD1), se non compensato, causa una deriva irreversibile dalla traiettoria (moto a spirale), rendendo indispensabile l'adozione di strategie avanzate di diagnosi e tolleranza.
    
    \item \textbf{Diagnosi dei Guasti (Cap. \ref{chap:fdi}):} Il banco di osservatori geometrici progettato ha mostrato una capacità di isolamento perfetta. I residui generati (\( r_1, r_2 \)) sono risultati strutturalmente disaccoppiati: un guasto sulla ruota destra ha attivato esclusivamente il monitor dedicato, senza generare falsi allarmi sull'altro canale.
    
    \item \textbf{Robustezza e Tolleranza (Cap. \ref{chap:ftc}):} Il sistema FTC attivo ha chiuso il cerchio. L'algoritmo di stima adattiva ha identificato correttamente la perdita di efficienza (\( \hat{w} \to w_{reale} \)) in meno di 5 secondi, e la riconfigurazione degli ingressi ha permesso al robot di mantenere la traiettoria desiderata senza alcuna deviazione visibile, anche a fronte di una perdita del 40\% di potenza su un motore.
\end{itemize}

\section{Considerazioni Finali}
Le simulazioni effettuate dimostrano che l'integrazione di un sistema FDI (Fault Detection and Isolation) con un controllo riconfigurabile (FTC) aumenta drasticamente l'affidabilità del sistema robotico.
Mentre il controllore classico fallisce in presenza di guasti asimmetrici, la soluzione proposta garantisce la continuità operativa ("mission completion") in autonomia.

Sviluppi futuri potrebbero includere l'estensione dell'analisi a guasti simultanei su entrambe le ruote o la validazione dell'algoritmo in presenza di rumore sui sensori (incertezza di misura). Tuttavia, allo stato attuale, il sistema progettato soddisfa pienamente le specifiche richieste per la gestione di guasti singoli.